% SEGMENT OLP-0030-S01
% Part:sets-functions-relations
% Chapter: size-of-sets
% Section: zig-zag

\documentclass[../../../include/open-logic-section]{subfiles}

\begin{document}

\olfileid{sfr}{siz}{zigzag}

\olsection{காண்டோரின் வளைவழி முறை}

\begin{explain}
சில “எளிய” பட்டியலாக்கங்களை ஏற்கெனவே கருதினோம்.
இப்போது சற்றுக் கடினமான ஒன்றைக் கருதுவோம்.
இயல் எண்களின் ஜோடிகளைக் கொண்ட கணத்தைக் கருதுக.
\oliflabeldef{sfr:set:pai:sec}{இதனை \olref[set][pai]{sec} இல்
பின்வருமாறு வரையறுத்தோம்:}{இதன் வரையறை:}
\[
\Nat \times \Nat = \Setabs{\tuple{n,m}}{n,m \in \Nat}
\]
இந்த வரிசைப்படுத்தப்பட்ட ஜோடிகளைப் பின்வருமாறு ஓர்
\emph{அட்டவணையில்} அமைக்கலாம்:
\[
\begin{array}{ c | c | c | c | c | c}
& \mathbf 0 & \mathbf 1 & \mathbf 2 & \mathbf 3 & \dots \\
\hline
\mathbf 0 & \tuple{0,0} & \tuple{0,1} & \tuple{0,2} & \tuple{0,3} & \dots \\
\hline
\mathbf 1 & \tuple{1,0} & \tuple{1,1} & \tuple{1,2} & \tuple{1,3} & \dots \\
\hline
\mathbf 2 & \tuple{2,0} & \tuple{2,1} & \tuple{2,2} & \tuple{2,3} & \dots \\
\hline
\mathbf 3 & \tuple{3,0} & \tuple{3,1} & \tuple{3,2} & \tuple{3,3} & \dots \\
\hline
\vdots & \vdots & \vdots & \vdots & \vdots & \ddots\\
\end{array}
\]
$\Nat \times \Nat$ இல் உள்ள ஒவ்வொரு வரிசைப்படுத்தப்பட்ட ஜோடியும்
அட்டவணையில் சரியாக ஒருமுறை இடம்பெறும் என்பது தெளிவு.
குறிப்பாக, $\tuple{n,m}$ என்பது $n$ ஆவது கிடைவரிசையிலும்
$m$ ஆவது நெடுவரிசையிலும் இடம்பெறும்.
ஆனால் இத்தகைய அட்டவணையின் உறுப்புகளை “ஒருபரிமாணப்”
பட்டியலாக எவ்வாறு அமைப்பது?
கீழுள்ள அட்டவணையின் முறை இதைச் செய்வதற்கான ஒரு வழியைக் காட்டுகிறது
(வேறு பல வழிகளும் உள்ளன):
\[
\begin{array}{ c | c | c | c | c | c | c}
& \mathbf 0 & \mathbf 1 & \mathbf 2 & \mathbf 3 & \mathbf 4 &\dots \\
\hline
\mathbf 0 & 0  & 1& 3 & 6& 10 &\ldots \\
\hline
\mathbf 1 &2 & 4& 7 & 11 & \dots &\ldots \\
\hline
\mathbf 2 & 5 & 8 & 12 & \ldots & \dots&\ldots \\
\hline
\mathbf 3 & 9 & 13 & \ldots & \ldots & \dots & \ldots \\
\hline
\mathbf 4 & 14 & \ldots & \ldots & \ldots & \dots & \ldots \\
\hline
\vdots & \vdots & \vdots & \vdots & \vdots&\ldots & \ddots\\
\end{array}
\]\noindent
இந்த முறை \emph{காண்டோரின் வளைவழி முறை} எனப்படுகிறது.
அது $\Nat \times \Nat$ ஐப் பின்வருமாறு பட்டியலாக்குகிறது:
\[
\tuple{0,0}, \tuple{0,1}, \tuple{1,0}, \tuple{0,2}, \tuple{1,1},
\tuple{2,0}, \tuple{0,3}, \tuple{1,2}, \tuple{2,1}, \tuple{3,0}, \dots
\]
இது பின்வரும் முடிவை நிறுவுகிறது:
\end{explain}

% SEGMENT OLP-0030-S02
\begin{prop}\ollabel{natsquaredenumerable}
$\Nat \times \Nat$ என்பது !!{enumerable}.
\end{prop}

% SEGMENT OLP-0030-S03
\begin{proof}
$f \colon \Nat \to \Nat\times\Nat$ என்பது,
ஒவ்வொரு $k \in \Nat$ ஐயும் காண்டோரின் வளைவழி அட்டவணையில்
$n$ ஆவது கிடைவரிசை, $m$ ஆவது நெடுவரிசை ஆகியவற்றின் மதிப்பாக
$k$ அமையும் $\tuple{n,m} \in \Nat \times \Nat$ உடன் இணைக்கட்டும்.
\end{proof}

% SEGMENT OLP-0030-S04
\begin{explain}
இந்த உத்தியை நேர்த்தியாகப் பொதுமைப்படுத்தவும் முடியும்.
எடுத்துக்காட்டாக, இயல் எண்களின் வரிசைப்படுத்தப்பட்ட மும்மடிகளைக்
கொண்ட பின்வரும் கணத்தைப் பட்டியலாக்க இதைப் பயன்படுத்தலாம்:
\[
\Nat \times \Nat \times \Nat = \Setabs{\tuple{n,m,k}}{n,m,k \in \Nat}
\]
$\Nat \times \Nat \times \Nat$ என்பதை,
$\Nat \times \Nat$ உடன் $\Nat$ இன் கார்டீசியன் பெருக்கமாகக்
கருதுகிறோம்; அதாவது,
\[
\Nat^3 = (\Nat \times \Nat) \times \Nat =
\Setabs{\tuple{\tuple{n,m},k}}{n, m, k
  \in \Nat }
\]
ஆகவே ஒரு அச்சை $\Nat$ இன் பட்டியலாக்கத்தாலும்,
மற்றொரு அச்சை $\Nat^2$ இன் பட்டியலாக்கத்தாலும் குறியிடுவதன் மூலம்,
ஓர் அட்டவணையைக் கொண்டு $\Nat^3$ ஐப் பட்டியலாக்கலாம்:
\[
\begin{array}{ c | c | c | c | c | c}
& \mathbf 0 & \mathbf 1 & \mathbf 2 & \mathbf 3 & \dots \\
\hline
\mathbf{\tuple{0,0}} & \tuple{0,0,0} & \tuple{0,0,1} & \tuple{0,0,2} & \tuple{0,0,3} & \dots \\
\hline
\mathbf{\tuple{0,1}} & \tuple{0,1,0} & \tuple{0,1,1} & \tuple{0,1,2} & \tuple{0,1,3} & \dots \\
\hline
\mathbf{\tuple{1,0}} & \tuple{1,0,0} & \tuple{1,0,1} & \tuple{1,0,2} & \tuple{1,0,3} & \dots \\
\hline
\mathbf{\tuple{0,2}} & \tuple{0,2,0} & \tuple{0,2,1} & \tuple{0,2,2} & \tuple{0,2,3} & \dots\\
\hline
\vdots & \vdots & \vdots & \vdots & \vdots & \ddots \\
\end{array}
\]
இவ்வாறு, காண்டோரின் வளைவழி முறை போன்ற ஒன்றைப் பயன்படுத்தி,
$\Nat^3$ இன் ஒரு பட்டியலாக்கத்தைப் பெறலாம்.
இதைத் தொடர்ந்து செய்து, எந்த இயல் எண் $n$ க்கும்
$\Nat^n$ இன் பட்டியலாக்கங்களைப் பெறலாம்.
எனவே பின்வரும் முடிவைப் பெறுகிறோம்:
\end{explain}

% SEGMENT OLP-0030-S05
\begin{prop}
ஒவ்வொரு $n \in \Nat$ க்கும் $\Nat^n$ என்பது !!{enumerable}.
\end{prop}

% SEGMENT OLP-0030-S06
\begin{prob}\label{sfr:siz:zigzag:prob:posint-n}
ஒவ்வொரு $n \in \Nat$ க்கும் $(\PosInt)^n$ என்பது
!!{enumerable} எனக் காட்டுக.
\end{prob}

% SEGMENT OLP-0030-S07
\begin{prob}\label{sfr:siz:zigzag:prob:posint-star}
$(\PosInt)^*$ என்பது !!{enumerable} எனக் காட்டுக.
\cref{sfr:siz:zigzag:prob:posint-n} ஐ நீங்கள் முன்கொள்ளலாம்.
\end{prob}

\end{document}
